\newpage
\thispagestyle{empty}
\mbox{}
\newpage

\addchapnonumber{C.\ \ Ma thèse en 180 secondes}\label{annex:MT180}


In this appendix is presented the text of my participation to the French contest \textit{Ma thèse en 180 secondes}. Both the original in French and the translated version are given down below.
The video recording with English subtitles is availbale on the Internet at the following adress: \url{https://www.youtube.com/watch?v=pGzU3e5XTz0}.

%\setlenght\parindent{0em}

    \subsection*{French (original)}
\begin{quote}    
    Un rayon lumineux traverse l’espace, court entre les astres,\\
    Lorsque la Terre sur son chemin l’intercepte.\\ 
    De l’énergie que ce rayon convoie, l’atmosphère a l’embarras du choix.\\ 
    Elle en reflète une partie et transmet le reste à la planète.\ 
    
    Sous ce flux de chaleur, et pour ne pas surchauffer,\\
    La Terre à l’équilibre se met à rayonner.\\ 
    Mais gaz et particules jouent au plafond de verre,\\ 
    Et s’ils sont en excédant c’est réchauffement par effet de serre.\
    
        C’est ça, en émettant du \ce{CO2} par exemple, on opacifie l’atmosphère, mais que dans un sens ! Les rayons du Soleil entrent comme avant, mais le rayonnement infrarouge que la Terre émet vers l’espace pour se refroidir, lui, est bloqué, et la chaleur confinée.
    Dans cet échange d’énergie les nuages, qui s’interposent, sont déterminants. 
    Prenez par exemple ce baigneur, sur la plage, qui sortant de l’eau sent les rayons du Soleil réchauffer sa peau. Soudain ! un cumulus — un gros nuage — passant par là projette son ombre sur le nageur ; il a froid !
    Au contraire, une fois la nuit tombée ce même nuage, parce qu’il piégera les infrarouges, aura un effet réchauffant.\
    
    Alors pour prédire le climat, on ne peut pas se passer de chercher à comprendre ces nuages. Et c’est tout l’objet de ma thèse : comment ils naissent, évoluent, précipitent ou s’évaporent. Pour finalement savoir s’ils réchauffent ou refroidissent.
    Le hic c’est  que les nuages ne se forment pas spontanément dans l’atmosphère...\
    
    Vous connaissez toutes et tous le miroir embué en sortant de la douche ? Eh bien les nuages c’est un petit peu pareil ; comme la buée ils ont besoin d’un support sur lequel se former. Dans l’atmosphère ils naissent sur de petites particules en suspension dans l’air, qu’on appelle aérosols.
    On en distingue deux types : les aérosols qui permettent la formation des gouttelettes nuageuses, qui composent entre autres le nuage de notre baigneur, et ceux qui forment les cristaux de glace des cirrus par exemple, vous savez ces hauts nuages filandreux qui ne font pas d’ombre.\
    
    En Arctique — la région que j’étudie — le réchauffement est 4 fois plus rapide qu’ailleurs. On y rencontre des nuages dits en phase mixte particulièrement importants pour le climat. Composés à la fois de gouttelettes et de cristaux de glace, ils sont très dépendants du type d’aérosols dans l’air. Mais  les modèles numériques de prévisions climatiques ne prennent pas en compte les aérosols, et donc représentent mal ces nuages.\
    
    Pour corriger ce problème, je traque les aérosols de l’Arctique afin de mieux comprendre leurs provenances et leur natures : seront-ils plus ou moins nombreux sur une planète qui se réchauffe ?  Cela me permet ensuite de développer — sur ordinateur — de tout nouveaux modèles de nuages. Je les teste et les ajuste en les comparant à de vraies observations, jusqu’à ce qu’enfin je parvienne à simuler des nuages qui naissent, évoluent et disparaissent comme ceux qu’on voit dans le ciel.\
    
    L’urgence du réchauffement climatique nous pousse à le comprendre, pour autant nos recherches ne peuvent pas nous dispenser d’engager des changements sociétaux profonds. Car ce n’est pas la planète qui est en danger, c’est nous ! Alors, donnons-nous la chance de nous sauver nous-mêmes.
\end{quote} 

    \subsection*{English (translated)}
    
\begin{quote}    
    A beam of light travels through space, racing among the stars,\\
    Until the Earth intercepts it on its path.\\
    The energy it conveys, the atmosphere can decide:\\
    Reflecting a part of it and transmitting the rest to the planet.\

    Under this flow of heat, and in order not to overheat,\\
    The Earth reaches equilibrium and begins to radiate.\\
    But gases and particles play a glass ceiling game,\\
    And if they are in excess,\\
    It is warming by greenhouse effect.\
    
    There you are !
    By emitting \ce{CO2} we make the atmosphere opaque,
    but only in one direction!
    The Sun's rays enter as before,
    but the infrared radiation that the Earth emits
    to cool itself is blocked, and the heat is trapped.
    In this energy exchange, clouds that interfere
    play a crucial role.
    For example, take a swimmer on the beach,
    as he emerges from the water he feels
    the Sun's rays warming his skin.
    Suddenly, a cumulus\----a big cloud\----passes by and casts its shadow on the swimmer; he gets cold!
    Conversely, once night falls,
    this same cloud because it will trap
    the infrared radiation, will have a warming effect.\
    
    So to predict the climate,
    we cannot do without understanding these clouds.
    And that is the focus of my thesis:
    how they form, evolve, precipitate, or evaporate.
    To ultimately determine if they warm or cool.
    The catch is that clouds do not form
    spontaneously in the atmosphere...\
    
    You all know the foggy mirror when you
    get out of the shower?
    Well, clouds are a bit like that; like fog,
    they need a surface on which to form.
    In the atmosphere, they form on small
    particles suspended in the air, called aerosols.
    They are of two types: aerosols that
    allow the formation of cloud droplets,
    which make up our bather's cloud, and
    those that form ice crystals in cirrus
    for example, you know, those high,
    filamentous clouds that don't cast shadows.\
    
    In the Arctic – the region I study –
    warming is four times faster than elsewhere.
    Mixed-phase clouds, particularly important
    for the climate, are found there.
    Composed of both droplets and ice crystals,
    they are highly dependent on the type of aerosols in the air.
    However, numerical climate prediction models
    do not account for aerosols and therefore poorly represent these clouds.\
    
    To tackle this problem, I track Arctic
    aerosols to better understand their origins and nature.
    Will they be more or less numerous on a
    planet that is warming?
    Next, this allows me to develop, on computer,
    new cloud models.
    I test and adjust them by comparing them to
    real observations until I finally manage to simulate
    clouds that form, evolve, and disappear like those we
    see in the sky.\
    
    The urgency of climate change drives us to
    understand it, but our researches cannot exempt us
    from making profound societal changes.
    It is not the planet that is in danger, it is us!
    So let's give us a chance to save ourselves.

\end{quote}    


